% 原书分章 PDF 第 24—28 页：第 25.7 节完整剂量—反应例子与第 25.8 节。
\section{例：因果剂量—反应曲线的单点估计}

关于因果剂量—反应曲线的估计量及其推断，参见 Kennedy 等（2015）。
Kennedy 等（2015）的估计量基于增广逆删失概率加权（augmented IPCW）估计量，
而非 CV-TMLE；所用平滑方法是更一般的局部多项式平滑，
而非我们采用的核平滑。

令 $O=(W,A,Y)\sim P_0$，其中 $Y\in\{0,1\}$，
$A\in[0,1]$ 连续取值。
令 $\mathcal M$ 为 $P_0$ 的统计模型，
它仅对给定 $W$ 时 $A$ 的条件密度 $g_0(a\mid W)$ 作出假设。
令 $\bar Q_0(A,W)=P_0(Y=1\mid A,W)$，
$Q_{W,0}$ 为 $W$ 的概率分布。
目标参数定义为
$\Psi(P)=E_PE_P(Y\mid A=a_0,W)$，其中 $a_0$ 为给定值。
令 $\Psi(P)(a)\equiv E_PE_P(Y\mid A=a,W)$。
对于带宽 $h$ 及 $J-1$ 阶正交核 $K$，令
$\Psi_h(P)=\int K_{h,a_0}(a)\Psi(P)(a)\,da$。
也可将该目标参数表示为 $\Psi_h(Q_W,\bar Q)$。
$\Psi_h$ 在 $P$ 处的有效影响曲线为
\[
\begin{aligned}
 D_h^*(P)
 &=\frac{K_{h,a_0}(A)}{g(A\mid W)}(Y-\bar Q(A,W))
 +\int_aK_{h,a_0}(a)\bar Q(a,W)\,da-\Psi_h(P)\\
 &\equiv D_h(\bar Q,g)-\Psi_h(P).
\end{aligned}
\]
也可将该有效影响曲线表示为 $D_h^*(\bar Q,g,\psi_h)$。有
\[
 \Psi_h(Q_W,\bar Q)-\Psi_h(Q_{W,0},\bar Q_0)
 =-P_0D_h^*(P)+R_{20,h}(P,P_0),
\]
其中
\[
 R_{20,h}(P,P_0)=E_{P_0}\int_a
 K_{h,a_0}(a)(\bar Q-\bar Q_0)(a,W)
 \frac{(g-g_0)(a\mid W)}{g(a\mid W)}\,da.
\]
给定交叉验证方案 $B_n\in\{0,1\}^n$，
令 $Q_{W,n,B_n},g_{n,B_n},\bar Q_{n,B_n}$
为将初始估计量应用于训练样本 $\{O_i:B_n(i)=0\}$ 得到的估计。
令
\[
 \mathrm{Logit}\,\bar Q_{n,B_n,\epsilon,h}
 =\mathrm{Logit}\,\bar Q_{n,B_n}+\epsilon C(h,a_0,g_{n,B_n}),
\]
其中 $C(h,a_0,g)(A,W)=K_{h,a_0}(A)/g(A\mid W)$。
令
\[
 \epsilon_n=\argmin_\epsilon
 E_{B_n}P^1_{n,B_n}L(\bar Q_{n,B_n,\epsilon,h}),
\]
其中
$L(\bar Q)(O)=-Y\log\bar Q(A,W)-(1-Y)\log(1-\bar Q(A,W))$
为 $\bar Q_0$ 的对数似然损失。
令 $\bar Q^*_{n,B_n,h}=\bar Q_{n,B_n,\epsilon_n,h}$。

现在还可以用一个最不利子模型对 $Q_{W,n,B_n}$ 作靶向更新，其得分为
\[
 D_{h,W}^*(Q_W,\bar Q)
 =\int_aK_{h,a_0}(a)\bar Q(a,W)\,da-\Psi_h(Q_W,\bar Q),
\]
从而有
\[
 E_{B_n}P^1_{n,B_n}
 D_{h,W}^*(Q^*_{W,n,B_n,h},\bar Q^*_{n,B_n,h})=0.
\]
例如，可以定义通用最不利模型
\[
 dQ_{W,n,B_n,\epsilon,h}=dQ_{W,n,B_n}
 \exp\left(\int_0^\epsilon
 D_{h,W}^*(Q_{W,n,B_n,x,h},\bar Q^*_{n,B_n,h})\,dx\right),
\]
并令
$\epsilon_n=\argmin_\epsilon E_{B_n}P^1_{n,B_n}L_1(Q_{W,n,B_n,\epsilon,h})$，
其中 $L(Q_W)=-\log q_W$ 为对数似然损失。
$Q_{W,n,B_n}$ 的靶向版本定义为
$Q^*_{W,n,B_n,h}=Q_{W,n,B_n,\epsilon_n,h}$。
此时的确有
\[
 E_{B_n}P^1_{n,B_n}
 D_{h,W}^*(Q^*_{W,n,B_n,h},\bar Q^*_{n,B_n,h})=0.
\]
于是，$\Psi_h(P_0)$ 的 CV-TMLE 定义为
$E_{B_n}\Psi_h(Q^*_{n,B_n,h})$，
其中 $Q^*_{n,B_n,h}=(Q^*_{W,n,B_n,h},\bar Q^*_{n,B_n,h})$。
不过，注意求解后一个方程意味着
\[
\begin{aligned}
 E_{B_n}\Psi_h(Q^*_{W,n,B_n,h},\bar Q^*_{n,B_n,h})
 &=E_{B_n}P^1_{n,B_n}\int_a
 K_{h,a_0}(a)\bar Q^*_{n,B_n,h}(a,\cdot)\,da.
\end{aligned}
\]
因此，对 $Q_{W,n,B_n}$ 的靶向更新，最终重新定义了目标参数的评价方式。
所以，对 $Q_{W,n,B_n}$ 作靶向更新，并将 CV-TMLE 定义为
$E_{B_n}\Psi_h(Q^*_{n,B_n,h})$，
等价于不对 $Q_{W,n,B_n}$ 作靶向更新，而直接把 CV-TMLE 定义为
\begin{equation}
 \psi^*_{n,h}=E_{B_n}P^1_{n,B_n}\int_a
 K_{h,a_0}(a)\bar Q^*_{n,B_n,h}(a,\cdot).
\tag{25.13}
\end{equation}
注意，
\[
 D_h^*(Q^*_{W,n,B_n,h},\bar Q^*_{n,B_n,h},g_{n,B_n})
 =D_h^*(\bar Q^*_{n,B_n,h},g_{n,B_n},\psi^*_{n,h}),
\]
因此 CV-TMLE 满足
\[
 E_{B_n}P^1_{n,B_n}
 D_h^*(\bar Q^*_{n,B_n,h},g_{n,B_n},\psi^*_{n,h})=0.
\]
将带宽选择器 $h_n$ 定义为下式的解：
\begin{equation}
 \frac{d}{dh_n}\psi_{n,h_n}
 +C_n\sigma_n\frac{d}{dh_n}(nh_n)^{-0.5}=0,
\tag{25.14}
\end{equation}
其中 $C_n$ 为使解存在的最小负常数或正常数。
例如，如有可能，选择 $C_n\in\{-1.96,1.96\}$。
这里 $\sigma_n^2$ 为下式的相合估计量：
\begin{equation}
 \sigma_0^2=\int K^2(y)\,dy\int_w
 \frac{\bar Q_0(1-\bar Q_0)(a_0,w)}{g_0(a_0\mid w)}\,dQ_{W,0}(w).
\tag{25.15}
\end{equation}
这里 $\psi_{n,h}$ 为 CV 一步估计量，
即 Kennedy 等（2015）中估计量的交叉验证版本：
\[
 \psi_{n,h}=E_{B_n}P^1_{n,B_n}
 D_h(\bar Q^*_{n,B_n,h},g_{n,B_n}).
\]
所提议的 $\psi_0$ 的 CV-TMLE 现在定义为
\[
 E_{B_n}\Psi_{h_n}(Q^*_{n,B_n,h_n})
 =E_{B_n}P^1_{n,B_n}\int_a
 K_{h_n,a_0}(a)\bar Q^*_{n,B_n,h_n}(a,\cdot)\,da.
\]

现在应用定理 25.1。
假定对某个 $\delta>0$，几乎处处有 $g_0(a_0\mid W)>\delta>0$。
利用 Cauchy--Schwarz 不等式，可以用
$h^{-0.5}\|\bar Q^*_{n,B_n,h_n}-\bar Q_0\|_{a_0,P_0}$
与 $h^{-0.5}\|g_{n,B_n}-g_0\|_{a_0,P_0}$ 之和，
控制
$\|D_h^*(Q^*_{n,B_n,h_n},g_{n,B_n})-D_h^*(Q_0,g_0)\|_{P_0}$，
其中对于某个固定且任意小的 $h_1>0$，
\[
 \|f\|_{a_0,P_0}^2
 =\int_w\int_yK^2(y)f(a_0+h_1y,w)^2\,dy\,dP_0(w).
\]
我们假定，对于某个多项式速率 $r(n)$，
\[
 \max\bigl(\|g_{n,B_n}-g_0\|_{a_0,P_0},
 \|\bar Q^*_{n,B_n,h_n}-\bar Q_0\|_{a_0,P_0}\bigr)=O_P(r(n)).
\]
还可知，当 $h\to0$ 时，$hP_0\{D_h^*(P_0)\}^2\to\sigma_0^2$。
剩下要证明二阶项条件：
\[
\begin{aligned}
 (nh)^{0.5}E_{P_0}\int_y
 &K(y)(\bar Q^*_{n,B_n,h_n}-\bar Q_0)(a_0+h_nya,W)\\
 &\quad\frac{(g_{n,B_n}-g_0)(a_0+h_ny\mid W)}
 {g_{n,B_n}(a_0+h_ny\mid W)}\,dy=o_P(1).
\end{aligned}
\]
利用 Cauchy--Schwarz 不等式，只需证明
\[
 (nh_n)^{0.5}
 \|\bar Q^*_{n,B_n,h_n}-\bar Q_0\|_{a_0,P_0}
 \|g_{n,B_n}-g_0\|_{a_0,P_0}=o_P(1).
\]
我们将假定它成立。现在可以应用定理 25.1，得到
\[
\begin{aligned}
 \sqrt{nh_n}(\psi^*_{h_n,n}-\psi_{h_n,0})
 &=\sqrt{nh_n}(P_n-P_0)D_{h_n}^*(P_0)+o_P(1)\\
 &\Rightarrow_dN(0,\sigma_0^2).
\end{aligned}
\]
假定 $a\mapsto\psi_0(a)=E_{P_0}\bar Q_0(a,W)$
在 $a=a_0$ 处 $J_0$ 次连续可微。则有
$h^{-J_0}(\psi_{h0}-\psi_0)\to B_0(J_0)$，其中
\begin{equation}
 B_0(J_0)\equiv\int_yK(y)\frac{y^{J_0}}{J_0!}\,dy\,\psi_0^{J_0}(a_0).
\tag{25.16}
\end{equation}
令 $J_0^*=\min(J_0,J)$。
应用定理 25.2 可证明，若选择最优速率
$h_n=n^{-1/(2J_0^*+1)}$，则
\[
 n^{-J_0^*/(2J_0^*+1)}
 (\psi^*_{h_n,n}-\psi_0)\Rightarrow_dN(B_0(J_0^*),\sigma_0^2).
\]
由于定理 25.2 的条件与定理 25.1 相同，
只是把 $K_{h,a_0}$ 替换为
$K_{1,h,a_0}=h\frac{d}{dh}K_{h,a_0}$，
而我们已证明后者是行为类似的核，因此条件的验证与上面相同。
应用定理 25.3 可证明，若 $C_n=O_P(1)$，则
$h_n\sim n^{-1/(2J_0^*+1)}$。
这证明了以下定理。

\begin{bookresult}{定理}{25.6}
考虑上面定义的 CV-TMLE $\psi^*_{h,n}$，见式 (25.13)，
以及式 (25.14) 的带宽选择器 $h_n$。
并考虑 $B_0(J)$ 的定义 (25.16)、式 (25.15) 的 $\sigma_0^2$，
以及对某个任意小的 $h_1>0$ 定义的
\[
 \|f\|_{a_0,P_0}
 =\sqrt{\int_w\int_yK^2(y)f(a_0+h_1y,w)\,dy\,dQ_{W,0}(w)}.
\]
作如下假设。
\begin{itemize}
\item 对某个 $\delta>0$，几乎处处有 $g_0(a_0\mid W)>\delta>0$。
\item 对某个未知的 $J_0$，
$a\mapsto\psi_0(a)=E_{P_0}\bar Q_0(a,W)$
在 $a_0$ 处 $J_0$ 次连续可微。
\item 对某个多项式速率 $r(n)$，
\[
 \max\bigl(\|g_{n,B_n}-g_0\|_{a_0,P_0},
 \|\bar Q^*_{n,B_n,h_n}-\bar Q_0\|_{a_0,P_0}\bigr)=O_P(r(n)).
\]
\item
\[
 (nh)^{0.5}
 \|\bar Q^*_{n,B_n,h_n}-\bar Q_0\|_{a_0,P_0}
 \|g_{n,B_n}-g_0\|_{a_0,P_0}=o_P(1).
\]
\item $C_n=O_P(1)$。
\end{itemize}
令 $J_0^*=\min(J_0,J)$。有
$h_n\sim n^{-1/(2J_0^*+1)}$，且
\[
 n^{-J_0^*/(2J_0^*+1)}
 (\psi^*_{h_n,n}-\psi_0)\Rightarrow_dN(B_0(J_0^*),\sigma_0^2).
\]
令 $h_n^*=h_n/\log n$ 为 $h_n$ 的欠平滑版本，
$\psi^*_{h_n^*,n}$ 为相应的 CV-TMLE。则
\[
 (nh_n^*)^{-0.5}(\psi^*_{h_n^*,n}-\psi_0)
 \Rightarrow N(0,\sigma_0^2).
\]
特别地，
\[
 \psi^*_{h_n^*,n}\pm1.96\sigma_n/(nh_n^*)^{0.5}
\]
是 $\psi_0$ 的渐近 $95\%$ 置信区间。
\end{bookresult}

\section{注记与延伸阅读}

Rosenblatt（1956）提出了核密度估计，Parzen（1962）证明了其渐近正态性。
Nadaraya（1964）与 Watson（1964）提出了核回归估计量。
Einmahl 和 Mason（2000）研究了当带宽即使在样本量增大时仍保持随机时，
核估计量的收敛性。
Einmahl 和 Mason（2005）得到了比本章所给结果更强的分布结果，
即允许带宽保持随机的核回归估计量的置信带。
关于核密度估计量研究进展的进一步综述，参见 Wied 和 Weißbach（2012）。

对于我们的因果例子，Kennedy 等（2015）此前研究过连续单时点处理下平均结局的估计。
Neugebauer 和 van der Laan（2007）将连续暴露的边际结构工作模型，
定义为真实剂量—反应曲线在有限维空间上的投影。
与本章讨论的参数不同，对该投影作索引的参数路径可微，
并且在一定正则条件下可以 $\sqrt n$ 速率估计。
Rosenblum 和 van der Laan（2010a）描述了对这些工作模型作索引的参数的 TMLE。
